% element: 10-s024:e01
% element: 10-s024:e02
% element: 10-s024:e03
\slajdid{10-s024}
\begin{frame}[label=10-s024]{Maksimalna struja i izbor kolektorskog otpornika}
\setlength{\abovedisplayskip}{4pt}\setlength{\belowdisplayskip}{4pt}\setlength{\abovedisplayshortskip}{3pt}\setlength{\belowdisplayshortskip}{3pt}
\[I_{O\max}=I_{OH}=-\frac{V_{CC}-V_{OH\min}}{R_C}\]
„Maksimum“ znači maksimalnu apsolutnu vrednost struje.\par
Ako je granica $I_{OH}=-20\,\mathrm{mA}$, struja $-25\,\mathrm{mA}$ nije dozvoljena.\par\bigskip\begin{columns}[T,onlytextwidth]\column{.47\textwidth}\textcolor{DEcrvena}{Manje $R_C$}\par\medskip
Veći strujni kapacitet izlaza.\column{.49\textwidth}\centering \textcolor{DEcrvena}{Veće $R_C$}\par\medskip
Veće pojačanje; uža prelazna zona.\end{columns}\bigskip\Rezultat{Kompromis pri izboru $R_C$.}
\end{frame}
